% ===================================================================
%  اللوحة رقم ١٠ — البحر. المرفأ.
%
%  Traduction, faite en 2026, en arabe levantin d'aujourd'hui, sur
%  l'ido de texto/io. Regles de la colonne : voir 10-lawha-01.tex.
%  Une seule scene, une seule numerotation. 37 blocs, 98 renvois.
%
%  DEUX MOTS DE CE TABLEAU SONT PARTIS D'ICI ET SONT REVENUS PAR
%  L'ITALIEN, ET ILS SONT TOUS LES DEUX DANS LE PORT DE GUERRE.
%  L'arsenal (21) est « الترسانة » : le mot est l'arabe دار الصناعة,
%  la maison de la fabrique, passe a Venise sous la forme arzanà —
%  celle que Dante emploie pour l'Arsenal — puis darsena, puis
%  arsenal dans toutes les langues d'Europe, et rentre en arabe sous
%  sa forme italienne. Et l'amiral (84) est « الأميرال », de أمير,
%  le commandant, devenu admiralis en latin medieval avant de revenir.
%  LE LIVRET FAIT ECRIRE A CETTE COLONNE DEUX MOTS QU'ELLE A PRETES
%  AU MONDE ET QU'ELLE RELIT EN TRADUCTION.
%
%  La colonne gujaratie avait eu કુલી au tableau 12 et શતરંજ au
%  tableau 13, a un tableau d'intervalle. Celle-ci en a deux dans le
%  meme tableau, ET LES DEUX SONT DE MARINE — ce qui est exactement
%  ce qu'on attend de la cote qui a appris a naviguer a la
%  Mediterranee. Byblos, Tyr, Sidon, Beyrouth, Tripoli : le port de ce
%  tableau n'a rien a decrire.
%
%  ET LE PHARE (93) EST « المنارة », QUI EST AUSSI LE MINARET. La
%  racine est ن-و-ر, la lumiere : une منارة est l'endroit d'ou vient
%  la clarte, que ce soit un feu pour les navires ou un appel pour les
%  fideles. Le francais, lui, dit « phare » d'apres l'ile de Pharos
%  devant Alexandrie. DEUX LANGUES, UNE TOUR, ET CHACUNE L'A NOMMEE
%  PAR CE QU'ELLE EN VOIT : l'une par la lumiere, l'autre par le lieu.
%
%  Pour le reste, ce tableau est a cette colonne ce que le tableau 10
%  etait a la gujaratie : la piece ou elle est chez elle. مرفأ le
%  port, باخرة le vapeur (3), أسطول la flotte, صاري le mat (51),
%  شراع la voile (13), مجداف la rame (47), دفّة le gouvernail (30),
%  مرساة l'ancre (77), بحّار le matelot, جمرك la douane (65), مدّ et
%  جزر le flux et le reflux, برزخ l'isthme (90), خليج la baie (95).
%  Ce qui a du etre transcrit tient sur une ligne, et c'est toujours
%  la meme chose : la MACHINE — الطوربيد, الفرقاطة, الترامواي.
% ===================================================================

%%K t10-apar-1 apar
\VUsaut{4.0mm}
\VUtitre{15.0pt}{60}{اللوحة رقم ١٠}
\VUsaut{3.0mm}
\VUpk{11.4pt}{40}{البحر. --- المرفأ.}
\VUsaut{3.0mm}

%%K t10-01-1 p
1. --- بالصيف، وقت بعض الناس عم يدوّروا عالهدوء والبرودة عالجبل، غيرهن
بيفضّلوا يروحوا عالساحل. هيك عملنا السنة الماضية، لأنه منحبّ
\VUgras{المحيط} \textsuperscript{(1)} كتير.

%%K t10-02-1 p
2. --- وأنا، بحسّ بمتعة كبيرة وأنا عم تأمّل هالامتداد الهايل من الميّ
اللي عم يمتدّ لحدّ \VUgras{الأفق} \textsuperscript{(2)}، وعم شوف
\VUgras{البواخر} \textsuperscript{(3)} الضخمة عم تسافر \VUgras{رحلة}
طويلة، اللي عم يركبوا عليها المسافرين الرايحين على أميركا.

%%K t10-03-1 p
3. --- لهيك أبي، اللي بيعرف ميلي، دايماً بيختار، لقضاية العطلة، شاطئ
كتير هادي، جنب واحد من \VUgras{مرافئنا الحربية} الكبيرة.

%%K t10-03-2 p
وبآخر إقامة إلنا، \VUgras{الأسطول} إجا رسا وراء \VUgras{كاسر الموج}
\textsuperscript{(4)} الضخم اللي بيسكّر \VUgras{المرفأ الحربي}
\textsuperscript{(5)} وبيحميه، وقدرت متل ما بدّي أتأمّل \VUgras{السفن
الحربية} المختلفة، \VUgras{الغوّاصة} \textsuperscript{(10)} الزغيرة اللي
بتغطس وبتختفي تحت الميّ، وكمان \VUgras{المدرّعة} \textsuperscript{(9)}
الضخمة اللي لهلّق ما كانت بتخاف غير من هجمات \VUgras{زورق الطوربيد}
\textsuperscript{(8)}.

%%K t10-tit-2 sub
\VUsaut{3.0mm}
\VUpk{10.2pt}{40}{المراكب الزغيرة.}
\VUsaut{3.0mm}

%%K t10-04-1 p
4. --- هالمراكب من زمان كانت بتعرف \VUgras{بمهارة} تلاقي النقطة
الضعيفة بالدرع المعدني التخين لتلزق فيها \VUgras{الطوربيد} الخطر، اللي
انفجاره بيغرّق المركب ؛ بس مع هيك كانت أقل خوفاً من الغوّاصة، لأنه
مدمّرات الطوربيد بتلحقها بسهولة.

%%K t10-05-1 p
5. --- وقدرت كمان شوف \VUgras{الطرّادات} \textsuperscript{(6)} السريعة
المدرّعة، اللي تقريباً منيعة متل المدرّعة الحقيقية، وكم \VUgras{سفينة
نقل} \textsuperscript{(12)}، وتلاتة أو أربع \VUgras{زوارق مدفعية}
\textsuperscript{(7)}، وأخيراً \VUgras{فرقاطة} \textsuperscript{(11)}.
\VUgras{ومنظر هالأسطول، لمّا بيناور ليرفع مراسيه، هوّي أكتر شي مشوّق.}

%%K t10-06-1 p
6. --- شو ما أكبر الفرق بالبنا بين هالكتل الفولاذ الكبيرة و\VUgras{المراكب
أبو تلات صواري} الأنيقة أو \VUgras{المراكب الشراعية}
\textsuperscript{(13)} الرشيقة اللي لهلّق بتنستعمل بالأسطول التجاري،
واللي شفتها عم تفوت عالمرفأ بكل شراعها، وأنا عم اتمشّى عالـ\VUgras{رصيف}
\textsuperscript{(14)}. والحقيقة إنّها بتخسر كتير من ميزاتها لمّا يكون
\VUgras{الهوا} معاكس. لازم تنزّل كل \VUgras{شراعها} لترجع تفوت عالحوض،
و\VUgras{القاطرة} \textsuperscript{(16)} بتصير ضرورية لتيجي تاخدها وتشدّها
لعند الرصيف، متل \VUgras{صنادل} \textsuperscript{(17)} عادية.

%%K t10-tit-3 sub
\VUsaut{3.0mm}
\VUpk{10.2pt}{40}{المرفأ التجاري.}
\VUsaut{3.0mm}

%%K t10-07-1 p
7. --- واللي بيشوّق بالمرفأ اللي عم احكي عنه إنّه ما فيه بس مرفأ
حربي، مرتّب بشكل صناعي بأشغال عملاقة، بس كمان \VUgras{مرفأ تجاري}
رائع بـ\VUgras{مصبّ} نهر كبير.

%%K t10-08-1 p
8. --- ولسان اليابسة اللي بيفصل الحوضين مُحصّن ومحمي بسلسلة
\VUgras{بطّاريات} و\VUgras{أبراج} عم تتقدّم عالبحر. وبدّه إذن خاص
ليتمشّى الواحد عالـ\VUgras{أسوار} \textsuperscript{(15)}. وأنا حصّلته
بسهولة، بواسطة عمّي، اللي هوّي مهندس \VUgras{أحواض البنا}
\textsuperscript{(18)}، ورحت زرت السفن اللي عم يبنوها تحت عنابر واسعة.

%%K t10-09-1 p
9. --- وحتى حضرت \VUgras{إنزال} مدرّعة، وبتذكّر اللحظة المؤثّرة اللي
حسّها كل الجمهور لمّا الكتلة الضخمة، متروكة لحالها، بلّشت تزحلق على
هيكلها المهدي وإجت عامت على ميّ النهر.

%%K t10-10-1 p
10. --- وللوصول لأحواض البنا، لازم يقطع الواحد \VUgras{ساحة التدريب}
\textsuperscript{(19)} وين سيّاف الحامية بيجوا كل يوم يتمرّنوا، وبيمرق
على \VUgras{جسر} \textsuperscript{(20)} زغير من حديد بيوصل ساحة العرض
بـ\VUgras{الترسانة} \textsuperscript{(21)}.

%%K t10-tit-4 sub
\VUsaut{3.0mm}
\VUpk{10.2pt}{40}{الترسانة والأحواض.}
\VUsaut{3.0mm}

%%K t10-11-1 p
11. --- ومع عمّي، زرت هالبنا الكبير المليان \VUgras{مدافع}
\textsuperscript{(22)} و\VUgras{قذائف} \textsuperscript{(23)} وقنابل. وشفت
كمان \VUgras{مصنع الحديد} \textsuperscript{(24)} اللي عم يتصنّع فيه
\VUgras{مراجل} \textsuperscript{(25)} البواخر.

%%K t10-12-1 p
12. --- وكتير قريب، في \VUgras{الأحواض} \textsuperscript{(26)} —
أحواض التصليح — و\VUgras{الأحواض العايمة} \textsuperscript{(27)} اللافتة
كتير، اللي فيها قدرت شوف عن قريب، على مركب عم ينصلّح،
\VUgras{الرفّاص} \textsuperscript{(28)}، و\VUgras{عارضة القاع}
\textsuperscript{(29)} و\VUgras{الدفّة} \textsuperscript{(30)}.

%%K t10-13-1 p
13. --- والمرفأ التجاري بيستاهل الدراسة متل المرفأ الحربي تماماً،
وقضّيت ساعات كتير عم بحملق عالأرصفة وين مربوطة \VUgras{بواخر الدواليب}
\textsuperscript{(31)} اللي بتطلع بالنهر، وعم بتفرّج على \VUgras{روافع
الصواري} \textsuperscript{(32)} وهيّي عم تشتغل، اللي عم ترفع حمول ضخمة متل
الريش.

%%K t10-tit-5 sub
\VUsaut{4.0mm}
\VUpk{10.2pt}{40}{المرشد.}
\VUsaut{3.0mm}

%%K t10-14-1 p
14. --- وأعجبت بـ\VUgras{بواخر عبر الأطلسي} \textsuperscript{(34)} وهيّي
عم تطلع من الخليج، و\VUgras{أعلامها} \textsuperscript{(35)} مفرودة
بالهوا، و\VUgras{مرشد} \textsuperscript{(36)} ماهر عم يقودها، اللي بيعرف
بشكل عجيب يمشي عالـ\VUgras{ممرّ} المعلّم بـ\VUgras{عوّامات}
\textsuperscript{(40)} وشواخص، وهوّي عم يتفادى \VUgras{المراكب
المسطّحة} \textsuperscript{(41)} اللي بتنقاد بصعوبة بشراعها المربّع
الوحيد. وميّزت عند \VUgras{المقدّم} \textsuperscript{(37)} --- راس المركب
--- \VUgras{مقصورات} \textsuperscript{(39)} الدرجة التانية، وعند
\VUgras{المؤخّر} \textsuperscript{(38)} --- كوثل المركب --- الصالونات
لركّاب الدرجة الأولى.

%%K t10-15-1 p
15. --- وبعض المرّات حضرت كمان تحميل \VUgras{صندل} \textsuperscript{(42)}
اللي هلّق كوّموا فيه بضايع \VUgras{المخزن} \textsuperscript{(43)}
و\VUgras{محطّة البحر} \textsuperscript{(44)}، اللي جابتها قطارات كتيرة من
\VUgras{خطّ الرصيف} \textsuperscript{(45)}.

%%K t10-tit-6 sub
\VUsaut{3.0mm}
\VUpk{10.2pt}{40}{متعة المجداف.}
\VUsaut{3.0mm}

%%K t10-16-1 p
16. --- كتير مرّات جدّفت بـ\VUgras{حوض المدّ} \textsuperscript{(53)}،
اللي بتلجأ عليه \VUgras{القوارب} \textsuperscript{(46)}، وكانت متعة إلي
حرّك \VUgras{المجداف} \textsuperscript{(47)}. وطلعت على \VUgras{سطح}
\textsuperscript{(48)} \VUgras{قارب شراعي} \textsuperscript{(49)}، وتسلّقت،
بـ\VUgras{الحبال} \textsuperscript{(52)}، عالـ\VUgras{صاري}
\textsuperscript{(51)} لحدّ \VUgras{العوارض} \textsuperscript{(50)}،
وبعدين رجّعت لفّتي \VUgras{بزورق خفيف} \textsuperscript{(54)} بين
\VUgras{مراكب الصيد} \textsuperscript{(55)} التقيلة، وين عم تنشف
\VUgras{الشباك} \textsuperscript{(56)}.

%%K t10-17-1 p
17. --- وعالـ\VUgras{رصيف} \textsuperscript{(58)} مرقت قدّامي أنواع
\VUgras{البضايع} \textsuperscript{(59)} المختلفة، بـ\VUgras{صناديق}
\textsuperscript{(60)} أو بـ\VUgras{بالات} \textsuperscript{(61)}. وكانوا
عم يعملوا \VUgras{حمولة} \textsuperscript{(63)} الصنادل، و\VUgras{روافع}
\textsuperscript{(62)} قوية عم تشيلها وتحطّها على \VUgras{كاميونات}
\textsuperscript{(64)} وقت سوّاق الكاميونات عم يصرّحوا عنها وعم يدفعوا
الرسوم بـ\VUgras{الجمرك} \textsuperscript{(65)}.

%%K t10-18-1 p
18. --- وأخيراً، لمّا وصلت عالـ\VUgras{جسر الدوّار} \textsuperscript{(66)}
أو عالـ\VUgras{هويس} \textsuperscript{(69)}، وأنا مارق قدّام
\VUgras{الكرّاكة} \textsuperscript{(68)} اللي عم تنضّف \VUgras{القناة}
\textsuperscript{(67)}، نزلت عالرصيف جنب \VUgras{عامود الإشارة}
\textsuperscript{(70)}.

%%K t10-19-1 p
19. --- وعالجهة التانية من هالرصيف بترسي \VUgras{الزوارق}
\textsuperscript{(71)} اللي بتجيب ضبّاط الأسطول عالبرّ. ومنظر يستاهل
تشوف \VUgras{البحّارة} وهنّي عم يرفعوا مجاديفهن بإيقاع، وقت القارب
المحمول عالـ\VUgras{موج} \textsuperscript{(72)} بيرسي بسهولة عند درج
النزول. ومن هالمحلّ، بيقدر الواحد يراقب أحسن \VUgras{المدّ والجزر}
وحركة \VUgras{المدّ} و\VUgras{الجزر} عالـ\VUgras{شاطئ}
\textsuperscript{(73)}.

%%K t10-tit-7 sub
\VUsaut{3.0mm}
\VUpk{10.2pt}{40}{النادي.}
\VUsaut{3.0mm}

%%K t10-20-1 p
20. --- وللرجعة عالبيت، أخدت \VUgras{الترامواي الكهربائي}
\textsuperscript{(74)} اللي \VUgras{موقفه} \textsuperscript{(75)} جنب
\VUgras{العامود} المحطوط قدّام نادي الضبّاط.

%%K t10-20-2 p
ومن \VUgras{شرفة} \textsuperscript{(76)} النادي، المزيّنة بـ\VUgras{مراسي}
\textsuperscript{(77)} ومدافع وقذائف، كتير مرّات شفت جمع رائع من ضبّاط
البحرية : \VUgras{ملازمين} \textsuperscript{(78)}، مع سيوفهن ؛
و\VUgras{نقباء} \VUgras{المدفعية} \textsuperscript{(79)}
و\VUgras{المشاة} \textsuperscript{(80)}، مع \VUgras{سيوفهن المعقوفة}.
ووراهن كان في \VUgras{بحّار} \textsuperscript{(81)} عم يناولهن
\VUgras{المنظار}، لمّا بدّهن يميّزوا شو عم يصير بعيد.

%%K t10-21-1 p
21. --- وشفت كمان \VUgras{ملازمين أوّل} \textsuperscript{(82)} شباب عم
ياخدوا التعليمات من \VUgras{قائدهن} \textsuperscript{(83)} أو من
\VUgras{الأميرال} \textsuperscript{(84)}، وقت \VUgras{رئيس عرفاء}
\textsuperscript{(85)} عم يستنّى ليحملها عالمركب.

%%K t10-22-1 p
22. --- ومن هالشرفة، المنظر رائع : بيطلّ الواحد على كل الشاطئ
بـ\VUgras{خيمه} \textsuperscript{(86)} و\VUgras{غرف التبديل}
\textsuperscript{(87)}، وبيشوف، بساعة الاستحمام، \VUgras{المستحمّين}
\textsuperscript{(88)} اللي عم يلعبوا بفرح قدّام \VUgras{الكازينو}
\textsuperscript{(89)}.

%%K t10-23-1 p
23. --- وبآخر الشاطئ، الساحل بيعمل \VUgras{شبه جزيرة}
\textsuperscript{(91)} \VUgras{صخرية} زغيرة، وين \VUgras{الموج الكبير}
\textsuperscript{(92)} بيجي ينكسر واللي عليها مبنيّة \VUgras{منارة}
\textsuperscript{(93)}، عم تدلّ بالليل البحّارة عالطريق. وهالشبه جزيرة
موصولة باليابسة بس بـ\VUgras{برزخ} \textsuperscript{(90)} كتير ضيّق.

%%K t10-tit-8 sub
\VUsaut{3.0mm}
\VUpk{10.2pt}{40}{منظر عام عالسواحل.}
\VUsaut{3.0mm}

%%K t10-24-1 p
24. --- و\VUgras{الجرف} \textsuperscript{(94)} عم يمتدّ، مشقّق من البحر
اللي عم يرسم فيه على كيفه سلسلة \VUgras{خلجان} \textsuperscript{(95)}
و\VUgras{رؤوس}، وعم يخلّيه بأحلى منظر.

%%K t10-24-2 p
وعالـ\VUgras{هضبة} \textsuperscript{(96)}، اللي بتطلّ على
\VUgras{المضيق} \textsuperscript{(98)}، في \VUgras{قلعة}
\textsuperscript{(97)} كتير مهمّة وكم « فيلّا ».

%%K t10-25-1 p
25. --- وهالمضيق بيفصل عن القارّة \VUgras{جزيرة} \textsuperscript{(99)}
كتير خطرة، محاطة بـ\VUgras{شعاب} \textsuperscript{(100)} وموج متكسّر، وين
بعض المرّات بتجنح قوارب وحتى سفن كبيرة. ومع سرعة النجدة ووصول
\VUgras{قوارب النجاة} السريع، هالـ\VUgras{غرق} مرعب، وبعواصف
الاعتدال، غرقت كم سفينة بناسها وبضايعها.

%%K t10-25-2 p
ومع كل هالجوار، \VUgras{المرفأ} كتير بيقصدوه البحّارة.

\VUsaut{6.0mm}
\VUfilet{22mm}
